\section{Normal log canonical indices and finite covers}\label{sec:index}

Two preliminary arguments will keep the induction within normal varieties.
The first recovers an lc index from the index of one normalized boundary
component and its Stein degree. The second separates finite actions on
the factors of a quasi-\'etale product cover. Throughout this section,
the inductive assertions $I_j,L_j$ are available for $j<n$; the assertion
$K_n$ is used only in Proposition~\ref{prop:lc-reduction}.

\subsection{Descent from a normalized boundary component}

We begin with an elementary descent observation. A rational
$q$-canonical form means a nonzero rational section of the $q$th tensor
power of the canonical line at the generic point. Its divisor is defined
on a normal variety by orders at prime divisors.

\begin{lemma}[Descent of a vertical principal divisor]\label{lem:vertical-descent}
Let $f:X\to Z$ be a contraction of normal projective varieties over a
field of characteristic zero. A rational function on $X$ whose divisor
has no horizontal component belongs to the subfield $k(Z)\subset k(X)$.
Suppose, in addition, that $(X,B)$ is a rational pair with
$K_X+B\sim_{\Q}0$, and that $q(K_{X_\eta}+B_\eta)\sim0$ on the generic
fibre. There are a rational $q$-canonical form $s$ and a rational Cartier
divisor $D\sim_{\Q}0$ on $Z$ such that
\begin{equation}\label{eq:index-vertical-form}
                 \Div(s)+qB=f^*D.
\end{equation}
\end{lemma}

\begin{proof}
The generic fibre $X_\eta$ is normal and projective, with
$H^0(X_\eta,\OO_{X_\eta})=k(Z)$. A rational function with vertical
divisor restricts to a function with neither zeros nor poles at any
prime divisor of $X_\eta$. Normality makes it regular, and hence a
member of $k(Z)$.

For the second assertion, extend a trivializing rational form on
$X_\eta$ to a rational $q$-canonical form $s$ on $X$, using any rational
top form on $Z$ to identify relative and absolute canonical lines.
The divisor $H=\Div(s)+qB$ is vertical. Choose $a>0$ and a rational
$aq$-canonical form $\tau$ with $\Div(\tau)+aqB=0$. The function
$s^{\otimes a}/\tau$ has divisor $aH$, so the first assertion writes
it as $f^*h$ for some $h\in k(Z)^*$. Thus
$H=f^*((1/a)\Div(h))$, as required.
\end{proof}

\begin{proposition}[Reduction to the zero-boundary klt index]\label{prop:lc-reduction}
Assume $L_j$ for every $j<n$ and $K_n$. Then $L_n$ holds: for every
finite set $\Gamma\subset[0,1]\cap\Q$ there is an integer
$c(n,\Gamma)>0$ such that
\[
                    c(n,\Gamma)(K_X+B)\sim0
\]
for every normal projective lc $n$-fold pair $(X,B)$ with coefficients
in $\Gamma$ and $K_X+B\sim_{\Q}0$.
\end{proposition}

\begin{proof}
The assumptions give the zero-boundary klt index bounds through
dimension $n$. Xu's klt induction theorem
\cite[Theorem~1.8]{Xu2019} therefore supplies the assertion when $(X,B)$
is klt. We prove the non-klt case without an slc index statement.

\emph{A Mori fibre space and a horizontal lc divisor.}
Take a crepant $\Q$-factorial dlt modification, and enlarge $\Gamma$
by $1$. Write $T=\lfloor B\rfloor\ne0$. For sufficiently small
rational $\epsilon>0$, the pair $(X,B-\epsilon T)$ is klt and
\[
                  K_X+B-\epsilon T\sim_{\Q}-\epsilon T
\]
is not pseudo-effective: its intersection with the $(n-1)$st power
of an ample divisor is negative. The klt minimal model program
\cite[Corollary~1.3.3]{BCHM2010} ends in a Mori fibre space
$f:X'\to Z$.

The full log canonical divisor remains crepant through this program.
Indeed, fix compatible canonical divisors and write
$r(K_X+B)=\Div(h)$. Divisorial contractions push forward this same
identity, and flips preserve it under their codimension-one
identification. On a common resolution, both pullbacks equal
$r^{-1}\Div(h)$. Thus the full pairs on all models are lc, even
though they need not be dlt. The exact order of their log canonical
divisors is unchanged. We henceforth write $(X,B)$ for the last
model.

If $\dim Z=0$, the klt log Fano pair $(X,B-\epsilon T)$ shows that
$X$ is of Fano type. The bounded complement theorem
\cite[Theorem~1.7]{Birkar2019}, with a fixed hyperstandard set containing
$\Gamma$, gives a bounded integer $c$ and a complement $B^+\ge B$
with $c(K_X+B^+)\sim0$. The effective divisor $B^+-B$ is numerically
trivial, hence zero by intersection with an ample class. This settles
that case.

Suppose $\dim Z>0$. On the Mori fibre space the surviving transform
of $T$ is relatively ample, since $\epsilon T$ is relatively
$\Q$-linearly equivalent to $-(K_X+B-\epsilon T)$. In particular,
some prime component $S$ of $T$ dominates $Z$. Let
$\nu:S^\nu\to S$ be its normalization. Divisorial adjunction gives
\begin{equation}\label{eq:index-adjunction}
             (K_X+B)|_{S^\nu}=K_{S^\nu}+\Theta.
\end{equation}
Here $(S^\nu,\Theta)$ is a normal lc pair and $\Theta$ is effective;
its coefficients belong to the rational DCC set of differents
determined by $\Gamma$ \cite{KollarMori1998}. Its adjoint is
$\Q$-linearly trivial. The lower-dimensional assertion $L_{n-1}$,
together with global ACC as in the conventions, therefore bounds a
trivializing multiple of $K_{S^\nu}+\Theta$.

\emph{A common form degree.}
The geometric generic fibre of $f$ is a normal projective lc pair
of dimension less than $n$, with the same coefficient set. The
assertions $L_j$, $j<n$, bound a trivializing multiple there, and
linear triviality descends in that same degree to $k(Z)$. To see the
last point, the corresponding divisorial sheaf becomes an invertible
trivial sheaf after algebraic closure. Invertibility descends
faithfully flatly. Flat base change for global sections then gives
a one-dimensional space over $k(Z)$, and its evaluation map is an
isomorphism because it becomes one after the field extension.

Choose a common integer $q$, depending only on $n$ and $\Gamma$,
which clears $\Gamma$ and trivializes both this fibre adjoint and
$K_{S^\nu}+\Theta$. Lemma~\ref{lem:vertical-descent} gives
\eqref{eq:index-vertical-form}, with $D\sim_{\Q}0$. At the generic
point of $S$, the form $s$ has precisely its logarithmic pole of
order $q$. Its $q$-fold residue is a nonzero rational
$q$-canonical form $v$ on $S^\nu$, and adjunction gives the equality
of actual divisors
\begin{equation}\label{eq:index-residue-divisor}
       \Div(v)+q\Theta=(f\circ\nu)^*D.
\end{equation}
For this equality one may first take a sufficiently divisible
tensor power and apply the usual pluricanonical adjunction
isomorphism; dividing the resulting equality of divisorial orders
gives \eqref{eq:index-residue-divisor}. No Cartier assumption in
degree $q$ along the whole of $S$ is needed.

\emph{Stein factorization and norm.}
Choose a trivializing $q$-log canonical form $v_0$ on
$(S^\nu,\Theta)$. Then $a=v/v_0\in k(S)^*$ satisfies
\[
                     \Div(a)=(f\circ\nu)^*D.
\]
Factor the map as $S^\nu\xrightarrow{g}Y\xrightarrow{h}Z$, where
$g$ is a contraction and $h$ is finite, with $Y$ normal. The first
part of Lemma~\ref{lem:vertical-descent} gives $a\in k(Y)^*$.
Injectivity of pullback for rational Cartier divisors under the
surjective map $g$ yields
\[
                          h^*D=\Div_Y(a).
\]
One can check this injectivity at a prime divisor of $Y$, where a
nonzero order remains nonzero at a prime divisor above it. The norm
identity for the finite morphism $h$ now gives
\begin{equation}\label{eq:index-norm}
       (\deg h)D=\Div_Z\bigl(N_{k(Y)/k(Z)}a\bigr).
\end{equation}

It remains to bound $\deg h$, not the degree of any gluing cover.
Apply the arithmetic Stein-degree theorem \cite[Theorem~1.1]{SD}
to $(X_\eta,B_\eta)$ over $k(Z)$, with its prime coefficient-one
divisor $S_\eta$ and parameter $t=1$. This fibre is normal,
integral and projective, and its field of global functions is
$k(Z)$. The relative algebraic closure of $k(Z)$ in $k(S_\eta)$
is exactly $k(Y)$. Hence $\deg h$ is bounded in terms of the fibre
dimension, and therefore in terms of $n$.

Equations~\eqref{eq:index-vertical-form} and \eqref{eq:index-norm}
show that $q(\deg h)(K_X+B)$ is principal. Taking the least common
multiple of the finitely many possible degrees, and then a common
multiple with the klt and base-point cases, proves the proposition.
The earlier crepant comparisons return the same multiple to the
original pair.
\end{proof}

The normality restriction in Proposition~\ref{prop:lc-reduction} is
important. The argument uses a single normalized component and takes
a field norm; it does not trivialize the log canonical sheaf on the
possibly reducible scheme $\lfloor B\rfloor$.

The same index hypothesis also controls the action of finite
automorphisms on log volume forms. The following consequence will be
used on components of degenerating fibres in
Section~\ref{sec:denominators}.

\begin{lemma}[Finite-order log volume characters]\label{lem:finite-character}
Assume $L_j$ in a fixed dimension $j$. There is an integer $c_j>0$
with the following property. Let $(V,E)$ be a normal projective lc
$j$-fold pair, with $E$ reduced, and let $\omega$ be a rational top
form satisfying $\Div(\omega)+E=0$. If $\gamma$ is a finite-order
automorphism of the pair and
$\gamma^*\omega=\lambda\omega$, then $\lambda^{c_j}=1$.
\end{lemma}

\begin{proof}
Form the quotient $\pi:V\to W=V/\langle\gamma\rangle$. Define
the boundary $E_W$ by the finite adjunction formula
\[
                       K_V+E=\pi^*(K_W+E_W).
\]
At a prime divisor with ramification index $e$, its coefficient is
$1-1/e$ if the upstairs coefficient is zero, and $1$ if the
upstairs coefficient is one. Thus the coefficients belong to the
fixed rational DCC set
$\{1-1/e:e\in\N\}\cup\{1\}$. The finite discrepancy formula
makes $(W,E_W)$ lc \cite[Proposition~5.20]{KollarMori1998}. A power of
$\omega$ invariant under the finite cyclic group descends to a
rational pluricanonical form downstairs; the ramification formula
shows that it trivializes a multiple of $K_W+E_W$. This also
establishes its $\Q$-Cartierness and $\Q$-linear triviality.

The assertion $L_j$ and global ACC supply a uniform $c_j$ for this
DCC set such that $c_j(K_W+E_W)\sim0$. Pull back a trivializing
log pluricanonical form. It is an invariant generator upstairs,
necessarily a scalar multiple of $\omega^{\otimes c_j}$, since
$V$ is projective and integral. Therefore $\lambda^{c_j}=1$.
For $j=0$ the same conclusion holds with $c_0=1$.
\end{proof}

\subsection{Quasi-\'etale covers and factor actions}

A finite map of normal varieties is \emph{quasi-\'etale} if it is
\'etale outside a set of codimension at least two. We write
$\Omega_V^{[p]}$ for the reflexive extension of $p$-forms from the
smooth locus. The elementary character criterion below records
that a bound on a volume character controls the exact global
canonical index, including Cartierness.

\begin{lemma}[Canonical index and the volume character]\label{lem:index-character}
Let $\pi:R\to X$ be a finite Galois quasi-\'etale cover of normal
projective varieties, with $K_R\sim0$. Choose a generating reflexive
top form $\omega$ on $R$, and write
$g^*\omega=\chi(g)\omega$ for the Galois action. Then
\[
                       mK_X\sim0\quad\Longleftrightarrow\quad\chi^m=1.
\]
\end{lemma}

\begin{proof}
If $\omega^{\otimes m}$ is invariant, it descends at the function
field level to a rational $m$-canonical form on $X$. Quasi-\'etaleness
means that its divisor has order zero at every prime divisor of
$X$. It therefore trivializes $\OO_X(mK_X)$ as a divisorial sheaf,
so $mK_X$ is principal. Conversely, a trivializing form downstairs
pulls back to a scalar multiple of $\omega^{\otimes m}$ and is
invariant.
\end{proof}

For a projective klt variety $X$ with $K_X\sim_{\Q}0$, the cyclic
cover attached to a primitive trivialization of its canonical
multiple is connected and quasi-\'etale. It has Cartier trivial
canonical divisor and is klt by finite adjunction, hence canonical:
its discrepancies are integers greater than $-1$. The singular
Beauville--Bogomolov decomposition theorem then supplies a finite
quasi-\'etale cover
\begin{equation}\label{eq:index-product-cover}
                 P=A\times\prod_{i=1}^r B_i\longrightarrow X.
\end{equation}
Here $A$ is abelian and each $B_i$ is an irreducible Calabi--Yau or
irreducible holomorphic symplectic variety in the singular sense
\cite[Definition~1.4 and Theorem~1.5]{HoringPeternell2019}. This theorem combines
the splitting and holonomy results of Druel and
Greb--Guenancia--Kebekus \cite{Druel2018,GGK2019}.

We use the following part of the definition of these nonabelian
blocks. They and all their connected normal finite quasi-\'etale
covers are canonical with trivial canonical divisor. Their
reflexive form algebras are generated by a top form in the
Calabi--Yau case, or by a generically nondegenerate two-form in
the symplectic case. In particular they have no reflexive
one-forms and no reflexive forms in degree one below their
dimension. Contraction with a volume form also gives the absence
of vector fields.

The cover $P\to X$ need not be Galois, and its Galois closure
need not remain a product. The next lemma is the comparison needed
to use the factors nonetheless. The argument follows the product
direction construction in \cite[Section~4.1]{U4}; we include the
regularity and semilinear details required here.

\begin{lemma}[Equivariant comparison with the factors]\label{lem:cover-comparison}
Let $k$ be a field of characteristic zero and let
$P=\prod_{i=1}^t B_i$ be a product of geometrically integral
projective factors whose geometric fibres are positive dimensional
blocks as above; include the abelian part, when present, as a single
factor. Let $R\to P$ be a geometrically integral normal finite
quasi-\'etale cover, and put $d=\dim R$. If a finite group $G$ acts
on $R$, then
there is a subgroup $G_0$ of index at most $d!$ and connected
finite quasi-\'etale covers $R_i\to B_i$ with the following
properties:
\begin{enumerate}
\item $R\to\prod_iR_i$ is finite and quasi-\'etale;
\item $G_0$ acts regularly on each $R_i$, and the projections
      $R\to R_i$ are equivariant;
\item a generating volume form on $R$ is proportional to the
      wedge product of the pulled-back volume generators on the $R_i$.
\end{enumerate}
The conclusions also hold for finite semilinear actions, interpreting
a factor transformation as an isomorphism to the corresponding
field-conjugate variety. In particular the $d!$th power of any
finite-order transformation preserves all factor fields.
\end{lemma}

\begin{proof}
\emph{Intrinsic preservation of the product directions.}
Pullback on the common smooth big open, followed by reflexive
extension, gives
\begin{equation}\label{eq:index-tangent-splitting}
                         \mathcal T_R=\bigoplus_{i=1}^t\mathcal E_i.
\end{equation}
There are no global homomorphisms between two distinct summands.
This can be checked after extending $k$ to an algebraic closure.
At least one of the corresponding factors is nonabelian,
say $B_i$. Fix general smooth points in all complementary factors.
The resulting slices of $R$ are normal and finite over $B_i$:
generic flatness and geometric normality in characteristic zero
allow this after shrinking the complementary parameter space.
The locus omitted from the quasi-\'etale map has codimension at
least two on a general slice. Each connected component of the
slice is consequently a connected normal quasi-\'etale cover
of $B_i$.

On the smooth big open of that component, $\mathcal E_i$ is its
tangent sheaf, while the complementary summands are trivial.
A homomorphism in either direction gives a vector field or a
one-form, which extends reflexively and vanishes by the block
properties. These slices cover a dense open of $R$, proving
the claimed vanishing.

The block projections in \eqref{eq:index-tangent-splitting} are
therefore central idempotents of the finite-dimensional algebra
$\operatorname{End}_R(\mathcal T_R)$. Its primitive central
idempotents give nonzero direct summands of positive generic
rank, so there are at most $d$ of them. The action of $G$
permutes this finite set. Its kernel $G_0$ has index at most
$d!$ and fixes each block projection, since a central idempotent
is a sum of primitive central idempotents. This reasoning is
unchanged for semilinear algebra automorphisms. In particular, the
idempotent argument is carried out over $k$ itself; it does not
require extending the finite group action to an algebraic closure.

\emph{Recovering the finite factor covers.}
Let $F_i$ be the relative algebraic closure of $k(B_i)$ in $k(R)$,
and let $R_i$ be the normalization of $B_i$ in $F_i$. Equivalently,
$R\to R_i\to B_i$ is the Stein factorization of the projection.
The intermediate fields $F_i$ are regular over $k$, because
$k(R)$ is regular over $k$; hence the $R_i$ are geometrically
integral. Stein factorization commutes here with algebraic closure:
the extensions $k(R)/F_i$ are regular in characteristic zero.
At the generic point, the complementary directions in
\eqref{eq:index-tangent-splitting} span
$\operatorname{Der}_{k(B_i)}k(R)$. For a finitely generated
extension in characteristic zero, the common constants of these
derivations are exactly the elements algebraic over $k(B_i)$.
Consequently $G_0$ preserves $F_i$ and induces birational factor
transformations.

The extension obtained by adjoining the other product factors
to $k(B_i)$ is regular, and hence is linearly disjoint from the
finite extension $F_i/k(B_i)$. Thus the function field of the
normal product $R_i\times\prod_{j\ne i}B_j$ is intermediate
between $k(P)$ and $k(R)$. Normalization gives finite maps
\[
          R\longrightarrow R_i\times\prod_{j\ne i}B_j
             \longrightarrow P.
\]
Ramification over a divisor of $B_i$ would persist in this
product and then at a prolongation to $k(R)$, contradicting
quasi-\'etaleness. Thus $R_i\to B_i$ is quasi-\'etale. The
product map $R\to\prod_iR_i$ is proper with finite fibres,
because its composite to $P$ is finite. Its image has dimension
$d$, equal to the dimension of the integral target product, so
it is surjective and finite. The same ramification argument
shows it is quasi-\'etale. In particular every fibre of
$R\to R_i$ has dimension $d-\dim R_i$.

\emph{Why the factor transformations are regular.}
Fix $g\in G_0$, and let $\Gamma_i$ be the closed graph of its
induced birational transformation on $R_i$. The two morphisms
from $R$ to $R_i$ give a surjection $R\to\Gamma_i$. All its
fibres have dimension at least $d-\dim R_i$. If the first
projection $\Gamma_i\to R_i$ had a positive-dimensional fibre,
its preimage in $R$ would have dimension strictly larger than
$d-\dim R_i$, contradicting the preceding equidimensionality.
Thus this projection is finite and birational. Normality of
$R_i$ makes it an isomorphism. Applying the same argument to
$g^{-1}$ proves that the induced transformation is an
automorphism. For a semilinear transformation, replace the
target by its field conjugate; the argument is identical.

Finally, each geometric $R_i$ is a quasi-\'etale cover of its
geometric block. It is canonical and has trivial canonical
divisor. These properties descend to $k$; linear triviality
descends by the same one-dimensional-section argument used in
Proposition~\ref{prop:lc-reduction}. The
product of its volume generators pulls back under the finite
quasi-\'etale product map to a generator of $K_R$. Any two
generators differ by a nonzero constant. This proves all
three assertions.
\end{proof}

The semilinear version applies to families over a curve as follows.
After a finite extension of the curve function field, define the
geometric product cover and its blocks. Take the normal Galois
closure of the resulting finite extension over the \emph{original}
total function field, and enlarge the curve field to its relative
algebraic closure in that closure. The latter is a finite Galois
extension, because the original total field is regular over the
original curve field. The closure is then geometrically integral
over its new constant field, and its finite Galois transformations
act semilinearly. On geometric fibres the comparison cover is
quasi-\'etale: over the original smooth locus purity makes the
initial covers \'etale, and the same is true of their conjugates
and Galois closure. After any further finite base extension the
same construction can be repeated. Lemma~\ref{lem:cover-comparison}
therefore separates the factor actions after taking a bounded
power, without requiring a bound on the degree of the comparison
cover.
